\documentclass{article}%
\usepackage[T1]{fontenc}%
\usepackage[utf8]{inputenc}%
\usepackage{lmodern}%
\usepackage{textcomp}%
\usepackage{lastpage}%
\usepackage{geometry}%
\geometry{margin=2cm}%
\usepackage{expex}%
\usepackage[T1]{fontenc}%
\usepackage[utf8]{inputenc}%
\usepackage[spanish]{babel}%
%
\title{Análisis Lingüístico: Gramatica Normativa Kaqchikel}%
\author{Generado por el TFM de Elías Carrasco}%
%
\begin{document}%
\normalsize%
\maketitle%
\section{Page 80}%
\label{sec:Page80}%
Ejemplo: %
\subsection{Uso de la consonante epentética /y/}%
\label{subsec:Usodelaconsonanteepenttica/y/}%

%
Cuando el sustantivo finaliza con vocal, se introduce la /y/ para %
separar la vocal del enclítico de persona. %
Ejemplos: %
\subsection{Techlal b'ib'aj toj mya' pakb'a'n yol / Pronombres personales para predicados no verbales}%
\label{subsec:Techlalbibajtojmyapakbanyol/Pronombrespersonalesparapredicadosnoverbales}%

%
De acuerdo al estudio de variación dialectal realizado por Pé %
rez, Jiménez y García (2000), se puede concluir que las formas %
de marcar el sujeto estativo son bastante regulares; por lo que 

%
\end{document}